% GENERATED FILE. Edit canonical lessons, then run npm run generate:books.
% canonical-source-hash: fnv1a64:97d1704d845f4ab3
% canonical-lessons: ZH-C368-zhishao, ZH-C368-maobing, ZH-C368-liaobuqi, ZH-C368-guli, ZH-C368-xuanze, ZH-R368-review-364-368

\chapter{Keep Going}
\label{ch:zh-keep-going}
\hlchaptermodality{368}

\begin{chapteropening}
\textbf{By the end of this chapter:} \emph{I can say at least, name a bad habit, encourage someone and help them choose.}

Complete the last lesson of chapter 368: I can say at least, name a bad habit, encourage someone and help them choose.
\end{chapteropening}

\section[zhìshǎo]{zhìshǎo — at least}

\label{lesson:ZH-C368-zhishao}



\begin{quote}
Before the new one: say the Mandarin for an insight, then the Mandarin for ridiculous.
\end{quote}

\subsection*{You'll want to know: zhìshǎo}
\emph{zhìshǎo} — \textquotedblleft{}at least\textquotedblright{}.

\begin{quote}
\emph{nǐ zhìshǎo yīnggāi zhǔnshí lái} — You should at least turn up on time.
\end{quote}

Its characters come later; for now it is a word for the ear.

\begin{cousinweb}
\emph{zhìshǎo} is built from \emph{zhì} and \zh{少} 'few', both inherited from Middle Chinese, \emph{tsyijH} and \emph{syewX} (Guangyun, in Baxter's transcription), which glosses the first 'to arrive' and the second 'not many'.
\end{cousinweb}

\begin{grammarlens}[title={What you've built}]
One more word for giving reasons.
\end{grammarlens}

\subsection*{Your turn}
\begin{itemize}
\raggedright
  \item \emph{Say it:} \emph{zhìshǎo}
  \item \emph{Say it:} \emph{zhìshǎo}, once more
  \item \emph{Say it:} \emph{nǐ zhìshǎo yīnggāi zhǔnshí lái} — You should at least turn up on time.
  \item \emph{Recall:} say the Mandarin for an insight, then the Mandarin for ridiculous, then say \emph{zhìshǎo} again
\end{itemize}

\begin{tcolorbox}[breakable,colback=teal!4,colframe=teal!35!black,title={Before you move on}]
What does zhìshǎo mean? (\textquotedblleft{}At least\textquotedblright{}.) Say it once more.
\end{tcolorbox}
\section[máobìng]{máobìng — fault; bad habit}

\label{lesson:ZH-C368-maobing}



\begin{quote}
Before the new one: say the Mandarin for ridiculous, then the Mandarin for at least.
\end{quote}

\subsection*{You'll want to know: máobìng}
\emph{máobìng} — \textquotedblleft{}fault; bad habit\textquotedblright{}.

\begin{quote}
\emph{tā yǒu yí ge xiǎo máobìng} — He has one little fault.
\end{quote}

Its characters come later; for now it is a word for the ear.

\begin{cousinweb}
\emph{máobìng} is built from \emph{máo} and \zh{病} 'an illness', both inherited from Middle Chinese, \emph{maw} and \emph{bjaengH} (Guangyun, in Baxter's transcription), which glosses the first, quoting an older dictionary, 'the hair of the brows and the like', and the second 'worry' and 'pain'. How the pair came to mean a fault is not settled.
\end{cousinweb}

\begin{grammarlens}[title={What you've built}]
One more word for giving reasons.
\end{grammarlens}

\subsection*{Your turn}
\begin{itemize}
\raggedright
  \item \emph{Say it:} \emph{máobìng}
  \item \emph{Say it:} \emph{máobìng}, once more
  \item \emph{Say it:} \emph{tā yǒu yí ge xiǎo máobìng} — He has one little fault.
  \item \emph{Recall:} say the Mandarin for ridiculous, then the Mandarin for at least, then say \emph{máobìng} again
\end{itemize}

\begin{tcolorbox}[breakable,colback=teal!4,colframe=teal!35!black,title={Before you move on}]
What does máobìng mean? (\textquotedblleft{}Fault; bad habit\textquotedblright{}.) Say it once more.
\end{tcolorbox}
\section[\texorpdfstring{\zh{了不起} (liǎobuqǐ)}{了不起 (liǎobuqǐ)}]{\zh{了不起} (liǎobuqǐ) — amazing; remarkable}

\label{lesson:ZH-C368-liaobuqi}



\begin{quote}
Before the new one: say the Mandarin for at least, then the Mandarin for a bad habit.
\end{quote}

\subsection*{You'll want to know: \zh{了不起}}
\textbf{\zh{了不起}} — \emph{liǎobuqǐ} — \textquotedblleft{}amazing; remarkable\textquotedblright{}.

\begin{quote}
\textbf{\zh{你真了不起}} \emph{nǐ zhēn liǎobuqǐ} — You're amazing.
\end{quote}

The tone of bù is lost here: between the two parts it is said light, bu.

\begin{cousinweb}
\textbf{\zh{了不起}} (\emph{liǎobuqǐ}) is built from \zh{了}, \zh{不} 'not' and \zh{起} 'to rise', all three inherited from Middle Chinese, \emph{lewX}, \emph{pjut} and \emph{khiX} (Guangyun, in Baxter's transcription), which glosses the first 'clever' and 'finished' and the third 'to rise'. How the three came to mean remarkable is not settled.
\end{cousinweb}

\begin{grammarlens}[title={What you've built}]
One more word for giving reasons.
\end{grammarlens}

\subsection*{Your turn}
\begin{itemize}
\raggedright
  \item \emph{Say it:} \emph{liǎobuqǐ}
  \item \emph{Say it:} \emph{liǎobuqǐ}, once more
  \item \emph{Say it:} \emph{nǐ zhēn liǎobuqǐ} — You're amazing.
  \item \emph{Recall:} say the Mandarin for at least, then the Mandarin for a bad habit, then say \emph{liǎobuqǐ} again
\end{itemize}

\begin{tcolorbox}[breakable,colback=teal!4,colframe=teal!35!black,title={Before you move on}]
What does \zh{了不起} mean? (\textquotedblleft{}Amazing; remarkable\textquotedblright{}.) Say it once more.
\end{tcolorbox}
\section[gǔlì]{gǔlì — to encourage}

\label{lesson:ZH-C368-guli}



\begin{quote}
Before the new one: say the Mandarin for a bad habit, then the Mandarin for remarkable.
\end{quote}

\subsection*{You'll want to know: gǔlì}
\emph{gǔlì} — \textquotedblleft{}to encourage\textquotedblright{}.

\begin{quote}
\emph{lǎoshī chángcháng gǔlì wǒmen} — The teacher often encourages us.
\end{quote}

Its characters come later; for now it is a word for the ear.

\begin{cousinweb}
\emph{gǔlì} is built from \emph{gǔ} 'a drum' and \emph{lì}, both inherited from Middle Chinese, \emph{kuX} and \emph{ljejH} (Guangyun, in Baxter's transcription), which glosses the first, quoting an older dictionary, 'to beat a drum', and the second 'to urge on'.
\end{cousinweb}

\begin{grammarlens}[title={What you've built}]
One more word for giving reasons.
\end{grammarlens}

\subsection*{Your turn}
\begin{itemize}
\raggedright
  \item \emph{Say it:} \emph{gǔlì}
  \item \emph{Say it:} \emph{gǔlì}, once more
  \item \emph{Say it:} \emph{lǎoshī chángcháng gǔlì wǒmen} — The teacher often encourages us.
  \item \emph{Recall:} say the Mandarin for a bad habit, then the Mandarin for remarkable, then say \emph{gǔlì} again
\end{itemize}

\begin{tcolorbox}[breakable,colback=teal!4,colframe=teal!35!black,title={Before you move on}]
What does gǔlì mean? (\textquotedblleft{}To encourage\textquotedblright{}.) Say it once more.
\end{tcolorbox}
\section[xuǎnzé]{xuǎnzé — to make a choice}

\label{lesson:ZH-C368-xuanze}



\begin{quote}
Before the new one: say the Mandarin for remarkable, then the Mandarin for to encourage.
\end{quote}

\subsection*{You'll want to know: xuǎnzé}
\emph{xuǎnzé} — \textquotedblleft{}to make a choice\textquotedblright{}.

\begin{quote}
\emph{wǒ juéde nǐ yīnggāi zìjǐ xuǎnzé} — I think you should make the choice yourself.
\end{quote}

Its characters come later; for now it is a word for the ear.

\begin{cousinweb}
\emph{xuǎnzé} is built from \emph{xuǎn} and \emph{zé}, both inherited from Middle Chinese, \emph{sjwenX} and \emph{draek} (Guangyun, in Baxter's transcription), which glosses each of them 'to choose'.
\end{cousinweb}

\begin{grammarlens}[title={What you've built}]
One more word for giving reasons.
\end{grammarlens}

\subsection*{Your turn}
\begin{itemize}
\raggedright
  \item \emph{Say it:} \emph{xuǎnzé}
  \item \emph{Say it:} \emph{xuǎnzé}, once more
  \item \emph{Say it:} \emph{wǒ juéde nǐ yīnggāi zìjǐ xuǎnzé} — I think you should make the choice yourself.
  \item \emph{Recall:} say the Mandarin for remarkable, then the Mandarin for to encourage, then say \emph{xuǎnzé} again
\end{itemize}

\begin{tcolorbox}[breakable,colback=teal!4,colframe=teal!35!black,title={Before you move on}]
What does xuǎnzé mean? (\textquotedblleft{}To make a choice\textquotedblright{}.) Say it once more.
\end{tcolorbox}
\section[(dialogue)]{Looking back: from Not Allowed to Keep Going}

\label{lesson:ZH-R368-review-364-368}



\begin{quote}
Say the Mandarin for to urge someone and to not allow.
\end{quote}

\subsection*{Your turn}
\begin{itemize}
\raggedright
  \item \emph{Say it:} to urge someone, to not allow, a standard to measure by, to give official approval, to remind someone
  \item \emph{Say it:} to suggest, but (in everyday speech), which goes to show, might as well, a word of advice
  \item \emph{Say it:} to avoid, to watch out, to keep at it, to give up, to be worth it
  \item \emph{Say it:} to get the wrong idea, to know in depth, to understand someone's reasons, an insight, ridiculous
  \item \emph{Say it:} at least, a bad habit, remarkable, to encourage, to make a choice
  \item \emph{Read it:} \zh{准} \zh{可} \zh{解} — each one aloud, then a word that holds it
\end{itemize}

\begin{tcolorbox}[breakable,colback=teal!4,colframe=teal!35!black,title={Before you move on}]
To urge someone? (\emph{quàn}.) To not allow? (\textbf{\zh{不准}}, \emph{bùzhǔn}.) A standard to measure by? (\emph{biāozhǔn}.) To give official approval? (\emph{pīzhǔn}.) To remind someone? (\emph{tíxǐng}.) To suggest? (\emph{jiànyì}.) But (in everyday speech)? (\textbf{\zh{可是}}, \emph{kěshì}.) Which goes to show? (\textbf{\zh{可见}}, \emph{kějiàn}.) Might as well? (\textbf{\zh{不如}}, \emph{bùrú}.) A word of advice? (\emph{zhōnggào}.) To avoid? (\emph{bìmiǎn}.) To watch out? (\textbf{\zh{当心}}, \emph{dāngxīn}.) To keep at it? (\emph{jiānchí}.) To give up? (\emph{fàngqì}.) To be worth it? (\emph{zhíde}.) To get the wrong idea? (\emph{wùhuì}.) To know in depth? (\textbf{\zh{了解}}, \emph{liǎojiě}.) To understand someone's reasons? (\textbf{\zh{理解}}, \emph{lǐjiě}.) An insight? (\textbf{\zh{见解}}, \emph{jiànjiě}.) Ridiculous? (\textbf{\zh{可笑}}, \emph{kěxiào}.) At least? (\emph{zhìshǎo}.) A bad habit? (\emph{máobìng}.) Remarkable? (\textbf{\zh{了不起}}, \emph{liǎobuqǐ}.) To encourage? (\emph{gǔlì}.) To make a choice? (\emph{xuǎnzé}.)
\end{tcolorbox}
